\clearpage
\originalchapter{独立附录: Judson四题选译与解答}
\label{ch:judson-independent}
\begin{center}
原著作者: Thomas W. Judson, Stephen F. Austin State University\\
中文选译、补证与独立题解: Codex (OpenAI助手)\\
修改版发布者: xhc144/mit-ocw-notes 项目\\
修改版日期: 2026年10月8日
\end{center}
本章是自足的独立作品, 从本版权页起至后附 GNU Free Documentation License 全文止, 使用 GFDL 1.2或任何后续正式版本. 本册前言至参考文献的原教学主体 (含第1--17章) 继续使用 CC BY-NC-SA 4.0. 依据 GFDL 第7节将两项独立作品汇编; 汇编不对本附录追加非商业条件, 也不把原教学主体改用 GFDL. 本附录题面与解答全部为中文, 后附英文只用于保留必要版权原文.

\noindent Copyright 1997 by Thomas W. Judson.\par
\noindent Copyright 1997--2013 by Thomas W. Judson.\par
\noindent Copyright 2026, xhc144/mit-ocw-notes project, for the Chinese selection, translation and independently written solutions.\par
\noindent Permission is granted to copy, distribute and/or modify this independent appendix under the terms of the GNU Free Documentation License, Version 1.2 or any later version published by the Free Software Foundation; with no Invariant Sections, no Front-Cover Texts, and no Back-Cover Texts. A copy of the license is included in the section entitled ``GNU Free Documentation License''.\par
以下另列两历史版一致的原著许可声明原文; 它指原著, 本修改版的许可范围由上段和本章标题页明确.\par
\noindent Permission is granted to copy, distribute and/or modify this document under the terms of the GNU Free Documentation License, Version 1.2 or any later version published by the Free Software Foundation; with no Invariant Sections, no Front-Cover Texts, and no Back-Cover Texts. A copy of the license is included in the appendix entitled ``GNU Free Documentation License''.\par
上列两项原著版权声明分别保留自2012、2013版. 本修改版同样没有不变章节、封面文字或封底文字. 修改版许可声明只指本独立附录; 附录允许依 GFDL 进行商业或非商业复制及改作. 原作者及 MIT 的名称不表示对中文稿的认可. 附录完整 LaTeX 源码与许可证随本册源码 ZIP 发布.

\par\addvspace{9pt}\Needspace{4\baselineskip}
\noindent{\bfseries\fontsize{12}{17}\selectfont 修订历史 (History)}\par
原著 \textit{Abstract Algebra: Theory and Applications}, Thomas W. Judson, Stephen F. Austin State University, 2012年8月11日版 (442页) 与2013年8月16日版 (444页), 由作者公开发布. 原件注明的网络位置为 \url{http://abstract.pugetsound.edu}; 作者现官方版本目录为 \url{https://judsonbooks.org/abstract-algebra-theory-and-applications/}, 公开源码位置为 \url{https://github.com/twjudson/aata}. 本次逐题核对的两份原 PDF 为 \url{https://judsonbooks.org/aata-files/aata-20120811.pdf} 与 \url{https://judsonbooks.org/aata-files/aata-20130816.pdf}.

2026年修改版《独立附录: Judson四题选译与解答》, 中文改作者 Codex (OpenAI助手), 发布者 xhc144/mit-ocw-notes 项目. 只选第3章第5、10、48、50题的跨版本稳定内容, 给出完整中文题面与独立解答, 对应 MIT 18.703 Spring 2013 作业1第3、4、10题及作业2第14题. 两版同号一致不证明教师当年使用的具体版次. 原第49题两版内容冲突, 保留未确定状态, 不把新版重复题或旧版候选冒充课程原题. 原英文书页、原书 TeX 及图像不嵌入本附录.

\par\addvspace{9pt}\Needspace{4\baselineskip}
\noindent{\bfseries\fontsize{12}{17}\selectfont 致谢 (Acknowledgements)}\par
以下保留原著致谢的内容及语气, 不是新增作品的背书. Judson 感谢下列审读者有益的评论与建议: David Anderson (University of Tennessee, Knoxville), Robert Beezer (University of Puget Sound), Myron Hood (California Polytechnic State University), Herbert Kasube (Bradley University), John Kurtzke (University of Portland), Inessa Levi (University of Louisville), Geoffrey Mason (University of California, Santa Cruz), Bruce Mericle (Mankato State University), Kimmo Rosenthal (Union College), Mark Teply (University of Wisconsin). 他还感谢 Steve Quigley、Marnie Pommett、Cathie Griffin、Kelle Karshick 以及 PWS 其他工作人员在本项目中的指导, 并表示与他们合作十分愉快. 原开源版本获得 National Science Foundation 支持 (Award 1020957).
\clearpage
